\chapter{Implementation and Reproduction Details}
\label{ch:appendix}

This appendix collects details that are needed to reproduce the results but did not fit in the
main text: the repository layout, the exact hyperparameters of the experiments, the reproduction
commands and the map of evidence artifacts.

\section{Repository Layout}
\label{sec:appendix-layout}

Table~\ref{tab:layout} shows the top-level layout of the code repository.

\begin{table}[ht]
\centering
\caption{Top-level layout of the project's code repository}
\label{tab:layout}
\small
\begin{tabular}{lL{8.6cm}}
\toprule
\textbf{Path} & \textbf{Contents} \\
\midrule
\texttt{src/thesis\_s2s/runtime/}  & cascade, full-duplex server, text-to-speech, session
                                     recording \\
\texttt{src/thesis\_s2s/bargein/}  & feature extraction, detector, playback controller,
                                     evaluation on recorded data \\
\texttt{src/thesis\_s2s/data/}     & dataset pipeline, speaker diarization, pair extraction,
                                     export and audit \\
\texttt{src/thesis\_s2s/eval/}     & measures, comparison tests and the fail-closed evidence
                                     aggregator \\
\texttt{src/thesis\_s2s/model/}    & the early experimental artifact, for historical
                                     reference only \\
\texttt{configs/}                  & configuration of the data, detector, measures and the
                                     seven direct-track generations \\
\texttt{scripts/}                  & preflight, training, re-evaluation, checkpoint selection
                                     and locked evaluations \\
\texttt{docs/}                     & predefined protocols, the review waiver document and
                                     recorded decisions \\
\texttt{results/}                  & write-once result artifacts \\
\texttt{tests/}                    & automated test suite \\
\texttt{third\_party/}             & lock file of upstream repositories and models \\
\bottomrule
\end{tabular}
\end{table}

\section{Hyperparameters of the Direct-Track Experiments}
\label{sec:appendix-hparams}

Table~\ref{tab:hparams} shows the exact hyperparameters of the seven experimental generations.
All runs used the AdamW optimizer and gradient checkpointing.

\begin{table}[ht]
\centering
\caption{Hyperparameters of the seven generations of low-rank adaptation in the direct track}
\label{tab:hparams}
\small
\begin{tabular}{lrrrrrl}
\toprule
\textbf{Version} & \textbf{Max} & \textbf{Window} & \textbf{Rank} &
\textbf{Seed} & \textbf{Peak} & \textbf{Trained} \\
 & \textbf{steps} & \textbf{(s)} & & & \textbf{memory (GB)} &
\textbf{embeddings} \\
\midrule
v1   & 8000 & 20  & 64  & 20260823 & 22.920   & upstream key \\
v2   & 2000 & 100 & 64  & 20260827 & \missing & upstream key \\
v3   & 500  & 100 & 128 & 20260827 & 26.504   & none \\
v4   & 500  & 100 & 128 & 20260827 & 28.640   & two text embeddings \\
v5   & 500  & 100 & 128 & 20260827 & 28.640   & two text embeddings \\
v6.2 & 200  & 12  & 64  & 20260901 & 16.577   & three text tensors \\
v7a  & 200  & 12  & 64  & 20260901 & 22.660   & three text tensors \\
\bottomrule
\end{tabular}
\end{table}

\claim{Versions v4 and v5 differ only in the first-codebook loss weight: 100 against 10. In
addition to the hyperparameters above, version v6.2 had scheduled text-input dropout from 0.25 to
0.75 and an audio loss weight of 0.1, and was trained and evaluated on the same 32 samples. Version
v7a was trained on the same 32 samples, with the inner-monologue text as a phoneme romanization and
word-level timing.}

\section{Hyperparameters of the Detector and the Responder}
\label{sec:appendix-detector}

Table~\ref{tab:hparams2} collects the configuration of the interruption detection module and the
text responder adapter.

\begin{table}[ht]
\centering
\caption{Configuration of the interruption detection module and the text responder adapter}
\label{tab:hparams2}
\small
\begin{tabular}{L{6.6cm}R{4.2cm}}
\toprule
\textbf{Parameter} & \textbf{Value} \\
\midrule
\multicolumn{2}{c}{\textbf{Interruption detection module}} \\
\midrule
Sampling rate & 16000 Hz \\
Frame length & 20 ms \\
Context window length & 400 ms \\
Number of MFCCs & 13 with deltas \\
Fundamental frequency search range & 60--400 Hz \\
Number of trees & 200 \\
Maximum tree depth & 4 \\
Learning rate & 0.08 \\
Minimum samples per leaf & 8 \\
Decision horizon after the acoustic onset & 120 ms \\
Selected decision threshold & 0.70 \\
Inference time per step on the CPU & 7.931 ms \\
\midrule
\multicolumn{2}{c}{\textbf{Text responder adapter}} \\
\midrule
Rank and scale factor & 16 and 32 \\
Dropout & 0.05 \\
Training epochs & 2 \\
Random seed & 20260906 \\
Eligible training rows & 2576 of 6419 \\
Eligible development rows & 12 of 131 \\
Response-token loss of the base model & 3.723 \\
Loss after epochs one and two & 2.643 and 2.562 \\
Reduction relative to the base model & 31.18\% \\
\bottomrule
\end{tabular}
\end{table}

\section{Reproduction Commands}
\label{sec:appendix-commands}

The following commands reproduce the main stages. The \texttt{PYTHONPATH} variable must point to the
\texttt{src} directory; the subcommand of each stage is listed in Table~\ref{tab:commands}.

\begin{table}[ht]
\centering
\caption{Reproduction commands of the project's main stages}
\label{tab:commands}
\small
\begin{tabular}{L{4.8cm}L{6.2cm}}
\toprule
\textbf{Stage} & \textbf{Subcommand} \\
\midrule
Corpus selection plan & \texttt{plan-conversation-corpus} \\
Window reconstruction & \texttt{prepare-conversation-episodes} \\
Audit of prepared windows & \texttt{audit-prepared-episodes} \\
Speaker diarization & \texttt{diarize-conversation-episodes} \\
Audit of diarization & \texttt{audit-diarized-episodes} \\
Noise-condition labels & \texttt{annotate-conversation-noise} \\
Corpus yield estimate & \texttt{estimate-conversation-yield} \\
Building response pairs & \texttt{build-conversations} \\
Audit of pairs & \texttt{audit-conversations} \\
Stereo training export & \texttt{export-moshi-data} \\
Audit of the export & \texttt{audit-moshi-data} \\
Running the system & \texttt{serve} \\
Aggregate evidence status & \texttt{evidence-status} \\
Hardware preflight & \texttt{gpu-preflight} \\
\bottomrule
\end{tabular}
\end{table}

The locked evaluations run as separate scripts in the \texttt{scripts} directory;
Table~\ref{tab:eval-scripts} lists them.

\begin{table}[ht]
\centering
\caption{Locked evaluation scripts}
\label{tab:eval-scripts}
\small
\begin{tabular}{L{4.6cm}L{6.4cm}}
\toprule
\textbf{Evaluation} & \textbf{Script} \\
\midrule
Mechanical panel of the cascade & \path{evaluate_cascade_validation.py} \\
Re-recognition measure & \path{evaluate_cascade_intelligibility.py} \\
Two-stage semantic test & \path{evaluate_qwen4b_responder_v2.py} \\
Interruption detector approximation & \path{evaluate_recorded_bargein_proxy.py} \\
Generating the figures of this report & \path{make_thesis_figures.py} \\
\bottomrule
\end{tabular}
\end{table}

\claim{These scripts operate on write-once artifacts and must not be used to reopen a final test or
revise a finalized selection. The figures of this report were generated from these artifacts with
\path{scripts/make_thesis_figures.py}.}

\section{Map of Evidence Artifacts}
\label{sec:appendix-evidence}

Table~\ref{tab:evidence-map} names, for each claim of this report, its reference artifact.

\begin{table}[ht]
\centering
\caption{Map of claims and reference artifacts. All paths are relative to the \texttt{results}
         directory.}
\label{tab:evidence-map}
\footnotesize
\begin{tabular}{L{4.6cm}L{7.4cm}}
\toprule
\textbf{Claim} & \textbf{Reference artifact} \\
\midrule
Aggregate status and acceptance conditions & \path{eval/EVIDENCE_STATUS.json} \\
Readable summary & \texttt{eval/SUMMARY.md} \\
Mechanical panel of the cascade & \path{eval/cascade_real_service_validation_panel.json} \\
Retrospective descriptive analysis & \path{eval/cascade_validation_descriptive_analysis.json} \\
Intelligibility approximation & \path{eval/cascade_intelligibility_proxy.json} \\
Semantic qualification stage & \path{eval/qwen4b_responder_v2_development_proxy.json} \\
Final semantic test & \path{eval/qwen4b_responder_v2_final_test_proxy.json} \\
Responder adaptation & \path{training/responder_lora_v1.json} \\
Interruption detector approximation & \path{eval/interrupt_recorded_proxy.json} \\
Synthetic detector test & \path{eval/interrupt_bench.json} \\
Dataset audit & \path{conversation_audit.json} \\
Training export report & \path{moshi_export_report.json} \\
Checkpoint selection of version one & \path{moshi_checkpoint_selection.json} \\
Held-out loss of version one & \path{moshi_heldout_model_eval.json} \\
Selection and runtime of version two & \path{moshi_v2_checkpoint_selection.json} \\
Result of version v6.2 & \path{moshi_v6_text_dropout_reevaluation.json} \\
Runtime of version v6.2 & \path{moshi_v6_text_dropout_runtime.json} \\
Server-side latency & \path{eval/latency_bench.json} \\
Human study & \path{eval/human_study.json} \\
Final repository audit & \path{release/final_audit.json} \\
\bottomrule
\end{tabular}
\end{table}

\section{Runtime Environment}
\label{sec:appendix-env}

Table~\ref{tab:env} shows the main components of the runtime environment and the identifiers of the
models used.

\begin{table}[ht]
\centering
\caption{Main components of the runtime environment and the models used}
\label{tab:env}
\small
\begin{tabular}{L{4.6cm}L{7.0cm}}
\toprule
\textbf{Component} & \textbf{Identifier} \\
\midrule
Text responder of the final system &
Qwen3-4B-Instruct-2507
\newline{\scriptsize cdbee75f17c01a7cc42f958dc650907174af0554} \\
Responder of the 9/9 mechanical panel and re-recognition &
Qwen2.5-0.5B-Instruct
\newline{\scriptsize 7ae557604adf67be50417f59c2c2f167def9a775} \\
Responder baseline in the semantic test & Qwen2.5-0.5B-Instruct \\
Base model of the direct track & Moshika, seven billion parameters \\
Audio codec of the direct track & Mimi \\
Speech recognition & hybrid RNNT-CTC model of the Conformer family \\
Text-to-speech & Piper with a Persian voice \\
Judge of the semantic test &
Qwen3.8-27B-FP8
\newline{\scriptsize 017b9c7af6b5689d5dd426a76e0bc077eb5ca20a} \\
Training and evaluation hardware & NVIDIA H100 NVL, 93.09 GB \\
\bottomrule
\end{tabular}
\end{table}

The exact revision and weight hash of each model are recorded in the lock file of upstream
repositories and repeated in each result artifact, so that the claim ``this is exactly the model that
was used'' is auditable.

\section{Building This Report}
\label{sec:appendix-build}

The Persian original of this report is typeset with XeLaTeX and the XePersian package; this English
edition uses XeLaTeX with polyglossia, so that the Persian examples in the text are typeset in the
font of the Persian edition. To reproduce both, running \texttt{./build.sh everything} in the report
directory is enough (\texttt{./build.sh en-all} for the English editions only); the script runs the
typesetting and bibliography chain in an isolated container, so no local TeX installation is needed.

Two technical points corrected while preparing the report may be useful to later users of the same
template: first, the name of the localization option of XePersian has changed in recent versions and
the previous option is no longer accepted; second, the right-to-left table patch of the bidi package
is not compatible with the current release of the array package, and its dated release must be
requested. Both corrections are recorded with comments in \texttt{styles/common.tex} (and the second
also in \texttt{styles/english.tex}, because polyglossia loads bidi for Persian).

A third point concerns the typesetting of numbers in the Persian edition. In right-to-left text, the
bidirectional algorithm splits a decimal number at the separator into two parts and reverses their
order; as a result the value 3.150 would be displayed as ``150.3''. For this reason, in the Persian
edition every measured quantity is placed in an explicit left-to-right run, while the structural
numbering (the numbers of chapters, sections, pages, tables and figures) keeps the automatic
Persian digits.
